%=====================================================================
%  《N=2 超对称规范理论中的 S-对偶导引》中文译注本 —— 导言区
%  底本：L. Alvarez-Gaumé, S. F. Hassan, arXiv:hep-th/9701069v1
%        (CERN-TH/96-371), Fortsch. Phys. 45 (1997) 159--236, 82 pp.
%  编译方式：xelatex main.tex（两遍，第二遍生成目录与交叉引用）
%
%  说明：
%  1) 本文档类为 ctexart（论文体例），章节编号与原文严格一致；
%  2) 公式环境沿用原文的 equation / eqnarray 结构，公式编号由 LaTeX
%     自动生成，与原文 (1.1)--(4.x) 逐一对应。原文的简写宏
%     \beq \eeq \bea \eea \bq \eq 已在整理底本时展开为
%     标准环境，以方便教授直接在源码上删改；
%  3) % [Tutor Note: ...] 为审校批注（LaTeX 注释，编译后不显示）；
%  4) 符号约定：原文全文采用 (+,-,-,-) 号差，本译本一律照录，
%     不逐式改写为 (-,+,+,+)（详见 front.tex 第 2 节“符号约定说明”）。
%=====================================================================
\usepackage[a4paper,left=2cm,right=2cm,top=2.4cm,bottom=2.4cm,
            headheight=14pt]{geometry}
\usepackage{amsmath}
\usepackage{amssymb}
\usepackage{amsthm}
\usepackage{mathtools}
\usepackage{bm}
\usepackage{graphicx}
\usepackage{xcolor}
\usepackage[most]{tcolorbox}
\usepackage{caption}
\usepackage{enumitem}
\usepackage{booktabs}
\usepackage{array}
\usepackage{multirow}
\usepackage{fancyhdr}
\usepackage[unicode,colorlinks=true,linkcolor=blue!45!black,
            citecolor=blue!45!black,urlcolor=blue!45!black]{hyperref}
\usepackage{bookmark}

%---------------------- 版式与字体 ----------------------
\allowdisplaybreaks[4]
\linespread{1.32}
\setlength{\emergencystretch}{3em}
\setlength{\parindent}{2em}
\setlist{itemsep=2pt,topsep=3pt,parsep=0pt}
\renewcommand{\figurename}{图}
\captionsetup{font=small,labelsep=quad,justification=centering,skip=5pt}
\ctexset{
  section={format=\Large\bfseries\heiti, beforeskip=18pt, afterskip=8pt},
  subsection={format=\large\bfseries, beforeskip=12pt, afterskip=5pt},
  subsubsection={format=\normalsize\bfseries, beforeskip=8pt, afterskip=3pt}
}

%---------------------- 页眉页脚 ----------------------
\pagestyle{fancy}
\fancyhf{}
\fancyhead[L]{\small $N=2$ 超对称规范理论中的 $S$-对偶导引}
\fancyhead[R]{\small\rightmark}
\fancyfoot[C]{\thepage}
\renewcommand{\headrulewidth}{0.4pt}
\fancypagestyle{plain}{\fancyhf{}\fancyfoot[C]{\thepage}%
  \renewcommand{\headrulewidth}{0pt}}
\renewcommand{\sectionmark}[1]{\markright{\thesection\quad #1}}

%---------------------- 本节重点（Key Takeaways） ----------------------
% 每个 \section 标题之后使用：
%   \begin{keypoints} \item ... \end{keypoints}
\newtcolorbox{keypoints}{
  enhanced, breakable,
  colback=blue!3!white, colframe=blue!55!black,
  boxrule=0.6pt, arc=2pt,
  left=8pt, right=8pt, top=4pt, bottom=4pt,
  title={本节重点（Key Takeaways）},
  fonttitle=\bfseries\small, fontupper=\small,
  coltitle=white, colbacktitle=blue!55!black
}

%---------------------- 常用记号（与原文一致，勿改） ----------------------
\newcommand{\rd}{\mathrm{d}}
\newcommand{\Tr}{\operatorname{Tr}}
\newcommand{\tr}{\operatorname{tr}}
\newcommand{\Imag}{\operatorname{Im}}
\newcommand{\Real}{\operatorname{Re}}


%---------------------- PDF 元数据 ----------------------
\hypersetup{
  pdftitle={N=2 超对称规范理论中的 S-对偶导引（中文译注本）},
  pdfauthor={L. Alvarez-Gaume, S. F. Hassan; 中文译注},
  pdfsubject={arXiv:hep-th/9701069v1, Fortschr. Phys. 45 (1997) 159-236},
  pdfkeywords={N=2 supersymmetry, Seiberg-Witten, S-duality, monopole, 中译本}
}
